\section{Cronograma de actividades (2026)}

Del 01 de junio al 21 de junio (Semanas 1 a 3).
\begin{itemize}
	\item Se utilizará la API de KEGG para la extracción automática de los datos correspondientes a las secuencias del grupo K01220.
	\item Se desarrollarán los scripts necesarios en Python implementados en notebooks de Jupyter.
	\item Se realizará la limpieza de los datos para generar de forma eficiente los archivos de salida en formato FASTA.
	\item Se registrará el procesamiento y análisis de los datos, alojando el código en el repositorio antes mencionado.
\end{itemize}

Del 22 de junio al 05 de julio (Semanas 4 y 5).
\begin{itemize}
	\item Se importarán las secuencias en formato FASTA al software MEGA11.
	\item Se realizará el alineamiento múltiple empleando el algoritmo MUSCLE.
	\item Se evaluará la conservación de regiones homólogas entre las secuencias.
	\item Se detectarán y documentarán las posibles secuencias divergentes dentro del grupo ortólogo analizado.
\end{itemize}

Del 27 de julio al 23 de agosto (Semanas 6 a 9).
\begin{itemize}
	\item Se procesarán las secuencias sin alineamiento previo utilizando la plataforma dbCAN.
	\item Se identificarán los clústeres de genes relacionados con carbohidratasas (CGCs) y se predecirán los substratos según la secuencia de aminoácidos.
	\item Se llevará a cabo la anotación de carbohidratasas utilizando perfiles de dominio derivados de la base de datos CAZy.
	\item Se ejecutarán tres estrategias de anotación: búsquedas basadas en modelos ocultos de Markov mediante pyHMMER, alineamientos de secuencia mediante DIAMOND y clasificación en subfamilias mediante dbCAN-sub.
	\item Se exportarán los resultados consolidados de dominios conservados y familias CAZy en archivos de Valores Separados por Tabulaciones.
\end{itemize}

Del 24 de agosto al 20 de septiembre (Semanas 10 a 13).
\begin{itemize}
	\item Se reconstruirán los árboles filogenéticos utilizando la herramienta IQ-TREE.
	\item Se aplicarán métodos de máxima verosimilitud para la estimación de soporte estadístico en las ramas.
	\item Se automatizará la selección de modelos evolutivos a través de la plataforma, lo que permitirá enfocar el análisis directamente en la lógica biológica de las relaciones evolutivas, delegando la notación y el modelado matemático complejo al software.
\end{itemize}

Del 21 de septiembre al 11 de octubre (Semanas 14 a 16).
\begin{itemize}
	\item Se utilizará la herramienta en línea DALI server para el análisis estructural de las secuencias de interés.
	\item Se basará el análisis en un alineamiento de matrices de distancia para buscar parentescos con estructuras proteicas previamente conocidas.
	\item Se cruzará la información con las bases de datos de AlphaFold Database, Protein Data Bank (PDB) y Pfam para consolidar la anotación de proteínas.
	\item Se escribirá el manuscrito de la tesina y se generarán las figuras pertinentes para la representación y presentación clara de los resultados obtenidos.
\end{itemize}